\documentclass{article}

% Track: workshop, double-blind. Add "final" only for the camera-ready
% version; omit it (and keep double-blind anonymization) for submission.
\usepackage[dblblindworkshop,final]{neurips_2026}

\usepackage[utf8]{inputenc}
\usepackage[T1]{fontenc}
\usepackage{hyperref}
\usepackage{url}
\usepackage{booktabs}
\usepackage{amsfonts}
\usepackage{amsmath}
\usepackage{amssymb}
\usepackage{nicefrac}
\usepackage{microtype}
\usepackage{xcolor}
\usepackage{graphicx}

\title{Working-Set-Aware Policy Gradients for Efficient Mixture-of-Experts Inference}
\workshoptitle{NeurIPS 2026 Workshop on Efficient Systems for Foundation Models}

\author{
  Diego Calanzone \\
  Mila, Quebec AI Institute \\
  \texttt{diego.calanzone@mila.quebec}
}

\begin{document}

\maketitle

\begin{abstract}
Serving very large mixture-of-experts (MoE) language models requires either tensor
parallelism across many accelerators or offloading expert weights to host memory, both
of which are constrained by the interconnect bandwidth available for expert transfers.
Prior work reduces this cost by encouraging the router to reuse a small set of experts
over a sequence, either through architectural changes to the routing mechanism, such as
option-critic style temporal extension, or through custom auxiliary losses that penalize
working-set misses while regularizing against distribution shift with a rank matching term. We
instead treat working-set locality as a reward signal and optimize it directly with an
established policy gradient algorithm, GRPO, using an LRU cache simulator over the
router's expert selections as a stateful reward function. This removes the need for
custom losses, additional trainable modules, or termination functions. We further combine
the working-set reward with the causal language modeling objective on the same rollouts,
which we interpret as jointly reinforcing the highest-reward completions sampled by the
policy and the ground truth completion, and find this to regularize distribution shift
more effectively than a KL penalty against a frozen reference. We introduce a prompt
conditioning variant that lets a single policy target different working-set budgets at
inference time by reading a literal instruction from its own context, and evaluate on
downstream task performance, working-set hit ratio, and stack distance, together with an
interpretability analysis of expert specialization.
\end{abstract}

\section{Introduction}

Deploying very large mixture-of-experts (MoE) language models, such as DeepSeek-V3 and
its successors \citep{liu2024deepseekv3} or the Kimi K2 line of models
\citep{moonshot2025kimi}, requires distributing hundreds of experts across multiple
accelerators, either through tensor parallelism or through offloading the bulk of expert
weights to host memory and streaming them to the device on demand
\citep{eliseev2023offloading}. Even with efficient serving systems, such as vLLM
\citep{kwon2023vllm}, throughput degrades sharply once expert transfers become frequent,
since the cost of moving weights over a comparably slow interconnect, e.g., PCIe, starts
to dominate over the cost of computation itself. This effect is worsened by the standard
load balancing objectives used at pretraining time, which explicitly encourage a router
to spread its selections uniformly over all experts to keep per-device utilization
even, at the expense of the locality that a working set would need to be useful at inference
time.

A natural way to alleviate this bottleneck is to fine-tune the router so that it
concentrates its choices on a smaller working set of experts per sequence, without
degrading the underlying model. One line of work draws on the option-critic framework
\citep{bacon2017optioncritic} and extends MoE routing temporally, holding an expert
selection as a macro-action until a learned termination function decides to switch
\citep{shen2026temporal}. This reduces the switch rate by design, but option-critic
architectures are known to be prone to degenerate solutions, such as options that never
terminate or that terminate at every step, and require an explicit deliberation cost to
avoid collapsing to one of these extremes \citep{harb2018waiting}, which introduces its
own training instability on top of the added termination and policy-over-options
modules. MELINOE \citep{raje2026melinoe} instead keeps the original routing architecture
and introduces a cache simulation loss that directly penalizes expert switching, paired
with a rank matching loss against a frozen reference router to prevent distributional
collapse onto a globally preferred subset of experts.

We take a different approach and ask whether working-set locality can be optimized
directly as a reward with an established policy gradient algorithm, rather than through
a bespoke loss or architectural change. We simulate a fixed-capacity LRU cache of
experts at a single router layer over the tokens a policy generates and use the
resulting hit ratio, or a soft relaxation of it, as the reward signal for GRPO and its
variant DAPO \citep{shao2024deepseekmath,yu2025dapo}, whose group-relative advantage and
trust region already provide most of the stability that a rank matching loss is meant
to supply, without an additional regularizer. We further combine this reward with the
causal language modeling objective on the model's own ground truth completions, which we
interpret as reinforcing two complementary trajectories per prompt at once, the
highest-reward rollout sampled by the policy and the one rollout known to be correct,
and find this to regularize distribution shift more effectively than a KL penalty
against a frozen reference or a distillation-style reward. Last, we introduce a prompt
conditioning variant that lets a single fine-tuned policy target different working-set budgets
at deployment time by reading a literal instruction from its own prompt.

Summarizing, we: i) reformulate working-set-friendly MoE fine-tuning as a policy gradient
problem with a stateful, simulator-based reward, requiring no custom loss or
architectural change; ii) show that combining this reward with a causal language
modeling objective on the same rollouts regularizes distribution shift more effectively
than a KL penalty; iii) introduce a prompt conditioning mechanism that lets one policy
serve multiple working-set budgets; and iv) evaluate the resulting policies on two MoE
backbones of different scale along task performance, working-set efficiency, and expert
specialization.

\section{Methodology}

\subsection{Stateful Working-Set Reward}
\label{sec:reward}

We simulate a fixed-capacity LRU cache of experts at a single transformer layer
$\ell$, chosen as the middle layer of the backbone. Let $p_t \in \Delta^{E}$ be the
router's distribution over $E$ experts at token $t$, and let $K$ be the number of
experts requested per token. We record a request set
$e_t = \mathrm{TopK}(p_t, K) \subset \{1, \dots, E\}$ at every position, including the
prompt, which is used only to warm the working set and is never rewarded. Let
$\mathcal{C}_t \subseteq \{1, \dots, E\}$, $|\mathcal{C}_t| \le c$, denote the set of
experts resident in the working set after token $t$, for a target capacity $c$. For each
$e \in e_t$, if $e \in \mathcal{C}_{t^-}$, the access is a hit and $e$ is promoted to
most-recently-used; otherwise it is a miss, $e$ is inserted, and the least-recently-used
expert is evicted once the capacity is exceeded. At training time, we route densely at
layer $\ell$, so that every expert participates in the forward pass weighted by $p_t$,
and use a soft reward, the probability mass already captured by the working set before it is
updated,
\begin{equation}
  r_t = \sum_{e \in \mathcal{C}_{t^-}} p_t[e], \qquad
  R(x, y) = \frac{1}{T} \sum_{t=1}^{T} r_t,
  \label{eq:softreward}
\end{equation}
where $T$ is the number of generated tokens and $R(x, y)$ is the sequence-level reward
handed to GRPO. This reward is stateful, since the value assigned to a token depends on
the entire working-set trajectory up to that point, but it requires no additional trainable
parameters and no assumption about which eviction policy the router should be tailored
to beyond the choice of $c$. We give the full derivation of the reward and its
relationship to the classical LRU stack algorithm in Appendix~\ref{app:stack}.

\subsection{Prompt Conditioning}
\label{sec:conditioning}

A single policy trained against a fixed working-set capacity $c$ only ever learns one
operating point. We instead let $c$ vary per example and condition the policy on it
through its own prompt, so that one fine-tuned model can be deployed at whichever working-set
budget the serving system actually has available. Given a pool of prompts
$\mathcal{D}$ and a set of candidate working-set sizes $\mathcal{C}$, we construct, for every
$x \in \mathcal{D}$ and every $c \in \mathcal{C}$, a conditioned prompt
$\tilde{x}_c = \texttt{[CACHE\_SIZE=}c\texttt{]} \oplus x$, prepending the literal
instruction to the first turn. The augmented training set
$\mathcal{D}' = \{(\tilde{x}_c, c) : x \in \mathcal{D}, c \in \mathcal{C}\}$ triples the
size of the original pool and is shuffled before training, so that every mini-batch
mixes examples conditioned on different budgets. Because the target capacity is now
encoded in the prompt itself, it is recovered at reward time by parsing the rollout's
own prompt text, $c = \mathrm{parse}(\tilde{x}_c)$, which makes $c$ a per-example
property of a training batch rather than a fixed hyperparameter, so the working set simulated
in Equation~\ref{eq:softreward} matches whichever budget the current example asks for.

\subsection{Combining Policy Gradients with Supervised Fine-Tuning}
\label{sec:sft}

GRPO estimates an advantage for a rollout $y_g$ by comparing its reward against the
mean and standard deviation of $G$ rollouts sampled from the same prompt,
$A_g = (R_g - \mathrm{mean}(R_{1:G})) / \mathrm{std}(R_{1:G})$, and the resulting policy
gradient reinforces whichever sampled rollout achieved the highest working-set reward relative
to its group. This can be read as a Monte Carlo estimate over the policy's own rollouts,
identifying and reinforcing the most working-set-friendly trajectory the policy currently
assigns non-negligible probability to. We additionally minimize the standard
teacher-forced negative log-likelihood of the dataset's own ground truth completion
$y^\star$ on the same prompt, so that the total objective is
\begin{equation}
  \mathcal{L}(\theta) = -\mathbb{E}_{x, \{y_g\}} \left[ \frac{1}{G} \sum_{g=1}^{G}
    A_g \log \pi_\theta(y_g \mid x) \right] + \lambda_{\mathrm{sft}}\, \mathrm{NLL}(y^\star \mid x),
  \label{eq:combinedloss}
\end{equation}
with no KL term against a frozen reference. Equation~\ref{eq:combinedloss} reinforces two
trajectories per prompt at once, the one rollout sampled by the policy that happened to
be the most working-set-friendly and the one rollout known a priori to be correct. When these
two objectives agree, i.e., when the ground truth completion is itself reasonably
working-set-friendly, which is common since natural language already reuses vocabulary and
concepts locally, their gradients concord; when they disagree, the supervised term keeps
the policy anchored to a distribution known to remain fluent and correct, which we find
to regularize distribution shift more effectively than a KL penalty or a
distillation-style reward, discussed further in Appendix~\ref{app:sft}.

\subsection{Metrics}
\label{sec:metrics}

\paragraph{Working-set hit ratio.} At evaluation time we replace the soft reward of
Section~\ref{sec:reward} with its hard counterpart, using deterministic top-$K$ routing
to match real deployment behavior rather than the differentiable, dense routing used
during training. The hit ratio of a rollout is
$h_t = \frac{1}{K} \sum_{e \in e_t} \mathbb{1}[e \in \mathcal{C}_{t^-}]$, averaged over
generated tokens.

\paragraph{Stack distance.} The hit ratio in Equation~\ref{eq:softreward} is a function
of the chosen capacity $c$, which makes it awkward to compare how working-set-friendly a
policy's routing is independently of any one deployment budget. We additionally report
the LRU stack distance \citep{mattson1970evaluation}: maintaining a recency-ordered
stack of distinct experts seen so far in a sequence, an access to expert $e$ at depth
$d$ in the stack, 1-indexed, has stack distance $d$, and is promoted to the top; an
access to an expert never seen before in the sequence is a compulsory miss. An access is
an LRU hit at capacity $c$ if and only if its stack distance is at most $c$, so a single
pass over a routing trace yields the hit ratio at every capacity at once, and the
distribution of stack distances itself, or a summary such as its mean, is a
capacity-independent measure of routing locality. We give the connection between this
quantity and the total number of expert transfers in Appendix~\ref{app:stack}.

\section{Experimental Setup}
\label{sec:setup}

We fine-tune two MoE backbones of different scale, Phi-tiny-MoE-instruct (32 layers, 16
experts, top-2) and OLMoE-1B-7B \citep{muennighoff2024olmoe} (16 layers, 64 experts,
top-8), using LoRA \citep{hu2021lora} adapters on the attention and expert projections.
Training data is drawn from the Nemotron Post-Training Dataset v2
\citep{nvidia2025nemotron}, filtered rather than truncated to fit the context budget. The
working-set reward is simulated at the middle layer of each backbone, with the working-set capacity
fixed as a fraction of the total number of experts and scaled accordingly across the two
backbones. We train with GRPO using the DAPO loss variant, group size $8$, and no KL
penalty when the supervised term of Section~\ref{sec:sft} is enabled. We compare five
objectives per backbone: causal language modeling alone; the working-set reward alone; their
combination, Equation~\ref{eq:combinedloss}; the prompt-conditioned variant of
Section~\ref{sec:conditioning}; and, as baselines, an option-critic style temporally
extended router \citep{shen2026temporal,bacon2017optioncritic} and MELINOE
\citep{raje2026melinoe}. Hyperparameters are given in Appendix~\ref{app:hparams}.

\section{Experimental Results}
\label{sec:results}

We report downstream task performance (MMLU \citep{hendrycks2021mmlu}, MMMLU, GSM8K
\citep{cobbe2021gsm8k}, HumanEval \citep{chen2021humaneval}, and MATH
\citep{hendrycks2021math}), working-set efficiency (hit ratio, stack distance, throughput), and an
interpretability analysis of expert specialization, comparing every objective of
Section~\ref{sec:setup} on both backbones. Results are organized by whether the policy
is conditioned on a target working-set size at inference time.

\paragraph{Note on LoRA.} Every objective below, including the base model's own baselines
and MELINOE and the option-critic controller, is evaluated as a LoRA adapter on top of the
frozen backbone rather than a fully fine-tuned model. Downstream task scores are
consequently expected to be somewhat lower across the board than a full fine-tune of the
same objective would reach, since LoRA's low-rank update is a strictly smaller hypothesis
class; this affects every method in this comparison equally and should not be read as a
property of any one objective. We also do not merge the adapter into the backbone before
timing throughput, so an adapter that touches the expert projections directly, as the
working-set reward does (Section~\ref{sec:setup}), adds its own inference-time compute
cost on top of whatever offload savings its routing behavior achieves; we control for this
in the throughput numbers below by timing the base model's own compute pass and adding
each objective's own working-set miss count on top, rather than each objective's own,
adapter-inflated, decode time.

\subsection{Without Prompt Conditioning}
\label{sec:results-unconditioned}

\begin{table}[h]
\centering
\small
\caption{Early, in-progress results on OLMoE-1B-7B (1024 held-out prompts). Phi-tiny-MoE
and the remaining objectives (causal LM alone, the combined reward of
Section~\ref{sec:sft}) are still running; this table is a placeholder to be replaced by
the full comparison. Working-set hit ratio and tokens/second are measured at the trained
capacity ($c=16$, top-8); tokens/second follows the estimated offloaded-inference latency
model of Appendix~\ref{app:stack} (compute plus a per-expert transfer penalty on every
working-set miss) rather than a wall-clock trace on real offloading hardware.}
\begin{tabular}{lccccccc}
\toprule
Objective & MMLU & MMMLU & GSM8K & HumanEval & MATH & Hit ratio & Tokens/s \\
\midrule
Base (untuned)         & 0.563 & 0.380 & 0.673 & 0.244 & 0.058 & 0.687 & 395.2 \\
Working-set reward      & 0.563 & 0.367 & 0.664 & 0.329 & 0.061 & \textbf{0.759} & 279.2 \\
Option-critic controller & 0.527 & 0.311 & 0.542 & 0.140 & 0.002 & 0.684 & 392.2 \\
MELINOE                & 0.492 & 0.322 & 0.435 & 0.079 & 0.002 & \textbf{0.763} & 439.5 \\
\bottomrule
\end{tabular}
\end{table}

The working-set reward alone (no supervised anchor) matches or improves the base model on
every task except MMMLU, and improves HumanEval by 8.5 points, while raising the hit
ratio by 7.2 points over the base model at the same capacity -- both option-critic and
MELINOE lose several points of downstream accuracy relative to the base model, most
sharply on MATH and GSM8K, for a comparable or smaller hit-ratio gain. Tokens/second is
computed against the base model's own compute pass plus each objective's own working-set
miss count (see the LoRA note above), which isolates the offload benefit from each
adapter's own inference-time cost; under this measure the working-set reward variant
still trails the base model rather than surpassing it, despite the higher hit ratio and
fewer raw misses, which we attribute to its shorter average completions (530 tokens
against the base model's 760) dividing a comparable total latency over fewer generated
tokens, and flag as a limitation of a purely token-averaged throughput metric rather than
evidence against the working-set benefit itself.

\subsection{With Prompt Conditioning}
\label{sec:results-conditioned}


\section{Conclusion}

We recast working-set-friendly fine-tuning of mixture-of-experts routers as a policy gradient
problem, using an LRU cache simulator as a stateful reward for GRPO in place of a custom
auxiliary loss or an architectural change to the routing mechanism. Combining this
reward with the causal language modeling objective on the same rollouts regularizes
distribution shift more effectively than a KL penalty against a frozen reference, and a
prompt conditioning variant lets a single policy serve multiple working-set budgets. We leave
the full experimental comparison, reported by working-set size and by whether the policy is
prompt-conditioned, to Section~\ref{sec:results}.

\clearpage
\bibliographystyle{plainnat}
\begin{thebibliography}{99}

\bibitem[DeepSeek-AI et al.(2024)]{liu2024deepseekv3}
DeepSeek-AI, Liu, A., Feng, B., Xue, B., et al.
\newblock DeepSeek-V3 Technical Report.
\newblock \emph{arXiv preprint arXiv:2412.19437}, 2024.

\bibitem[Moonshot AI(2025)]{moonshot2025kimi}
Moonshot AI.
\newblock Kimi K2: Open Agentic Intelligence.
\newblock \emph{Technical Report}, 2025.

\bibitem[Kwon et al.(2023)]{kwon2023vllm}
Kwon, W., Li, Z., Zhuang, S., Sheng, Y., Zheng, L., Yu, C.~H., Gonzalez, J.~E., Zhang, H., and Stoica, I.
\newblock Efficient Memory Management for Large Language Model Serving with PagedAttention.
\newblock \emph{Proceedings of the 29th Symposium on Operating Systems Principles}, 2023.

\bibitem[Eliseev and Mazur(2023)]{eliseev2023offloading}
Eliseev, A. and Mazur, D.
\newblock Fast Inference of Mixture-of-Experts Language Models with Offloading.
\newblock \emph{arXiv preprint arXiv:2312.17238}, 2023.

\bibitem[Shen and Henderson(2026)]{shen2026temporal}
Shen, Y. and Henderson, P.
\newblock Temporally Extended Mixture-of-Experts via Option-Critic.
\newblock \emph{arXiv preprint}, 2026.

\bibitem[Bacon et al.(2017)]{bacon2017optioncritic}
Bacon, P.-L., Harb, J., and Precup, D.
\newblock The Option-Critic Architecture.
\newblock \emph{Proceedings of the AAAI Conference on Artificial Intelligence}, 2017.

\bibitem[Harb et al.(2018)]{harb2018waiting}
Harb, J., Bacon, P.-L., Klissarov, M., and Precup, D.
\newblock When Waiting is not an Option: Learning Options with a Deliberation Cost.
\newblock \emph{Proceedings of the AAAI Conference on Artificial Intelligence}, 2018.

\bibitem[Raje et al.(2026)]{raje2026melinoe}
Raje, A., Nayak, A., and Joshi, A.
\newblock MELINOE: Fine-Tuning Mixture-of-Experts Language Models for Cache-Friendly Offloaded Inference.
\newblock \emph{arXiv preprint arXiv:2602.11192}, 2026.

\bibitem[Shao et al.(2024)]{shao2024deepseekmath}
Shao, Z., Wang, P., Zhu, Q., Xu, R., Song, J., Bi, X., Zhang, H., Zhang, M., Li, Y.~K., Wu, Y., and Guo, D.
\newblock DeepSeekMath: Pushing the Limits of Mathematical Reasoning in Open Language Models.
\newblock \emph{arXiv preprint arXiv:2402.03300}, 2024.

\bibitem[Yu et al.(2025)]{yu2025dapo}
Yu, Q., Zhang, Z., Zhu, R., Yuan, Y., Zuo, X., Yue, Y., Fan, T., Liu, G., Liu, L., Liu, X., et al.
\newblock DAPO: An Open-Source LLM Reinforcement Learning System at Scale.
\newblock \emph{arXiv preprint arXiv:2503.14476}, 2025.

\bibitem[Mattson et al.(1970)]{mattson1970evaluation}
Mattson, R.~L., Gecsei, J., Slutz, D.~R., and Traiger, I.~L.
\newblock Evaluation Techniques for Storage Hierarchies.
\newblock \emph{IBM Systems Journal}, 9(2):78--117, 1970.

\bibitem[Muennighoff et al.(2024)]{muennighoff2024olmoe}
Muennighoff, N., Soldaini, L., Groeneveld, D., Lo, K., Morrison, J., Min, S., Shi, W., Walsh, P., Tafjord, O., Lambert, N., et al.
\newblock OLMoE: Open Mixture-of-Experts Language Models.
\newblock \emph{arXiv preprint arXiv:2409.02060}, 2024.

\bibitem[Hu et al.(2021)]{hu2021lora}
Hu, E.~J., Shen, Y., Wallis, P., Allen-Zhu, Z., Li, Y., Wang, S., Wang, L., and Chen, W.
\newblock LoRA: Low-Rank Adaptation of Large Language Models.
\newblock \emph{arXiv preprint arXiv:2106.09685}, 2021.

\bibitem[Hendrycks et al.(2021)]{hendrycks2021mmlu}
Hendrycks, D., Burns, C., Basart, S., Zou, A., Mazeika, M., Song, D., and Steinhardt, J.
\newblock Measuring Massive Multitask Language Understanding.
\newblock \emph{Proceedings of the International Conference on Learning Representations}, 2021.

\bibitem[Cobbe et al.(2021)]{cobbe2021gsm8k}
Cobbe, K., Kosaraju, V., Bavarian, M., Chen, M., Jun, H., Kaiser, L., Plappert, M., Tworek, J., Hilton, J., Nakano, R., et al.
\newblock Training Verifiers to Solve Math Word Problems.
\newblock \emph{arXiv preprint arXiv:2110.14168}, 2021.

\bibitem[Chen et al.(2021)]{chen2021humaneval}
Chen, M., Tworek, J., Jun, H., Yuan, Q., Pinto, H.~P.~d.~O., Kaplan, J., Edwards, H., Burda, Y., Joseph, N., Brockman, G., et al.
\newblock Evaluating Large Language Models Trained on Code.
\newblock \emph{arXiv preprint arXiv:2107.03374}, 2021.

\bibitem[Hendrycks et al.(2021)]{hendrycks2021math}
Hendrycks, D., Burns, C., Kadavath, S., Arora, A., Basart, S., Tang, E., Song, D., and Steinhardt, J.
\newblock Measuring Mathematical Problem Solving with the MATH Dataset.
\newblock \emph{Proceedings of the Neural Information Processing Systems Track on Datasets and Benchmarks}, 2021.

\bibitem[NVIDIA(2025)]{nvidia2025nemotron}
NVIDIA.
\newblock Nemotron Post-Training Dataset v2.
\newblock \emph{Hugging Face Dataset Card}, 2025.

\bibitem[Jiang et al.(2024)]{jiang2024mixtral}
Jiang, A.~Q., Sablayrolles, A., Roux, A., Mensch, A., Savary, B., Bamford, C., Chaplot, D.~S., de las Casas, D., Hanna, E.~B., Bressand, F., et al.
\newblock Mixtral of Experts.
\newblock \emph{arXiv preprint arXiv:2401.04088}, 2024.

\end{thebibliography}

\appendix

\section{Stack Distance and the Number of Expert Transfers}
\label{app:stack}

We make precise the relationship between the stack distance defined in
Section~\ref{sec:metrics} and the total number of CPU-GPU expert transfers an offloaded
serving system would incur. Let $r^{(\ell,t)} \in \{0,1\}^E$ denote the binary request
vector at layer $\ell$ and token $t$, with $\|r^{(\ell,t)}\|_1 = K$, and let
$\mathcal{C}^{(\ell,t)}$ be the set of experts resident in the working set after token $t$
under an LRU eviction policy of capacity $c$. The number of working-set misses, which equals
the number of transfers under a cold-start-free serving system, is
\begin{equation}
  N_{\mathrm{miss}}(\ell, c) = \sum_{t=1}^{T} \sum_{e \in e_t} \mathbb{1}\!\left[
    e \notin \mathcal{C}^{(\ell,t^-)}\right] = \sum_{t=1}^{T} \sum_{e \in e_t}
    \mathbb{1}[d(e, t) > c],
\end{equation}
where $d(e, t)$ is the stack distance of the access to $e$ at token $t$, treating a
compulsory miss as $d = \infty$. Since $d(e, t)$ does not depend on $c$, computing it
once for every access yields $N_{\mathrm{miss}}(\ell, c)$ for every candidate capacity
$c$ in a single pass over the routing trace, which is how we obtain the hit ratio curves
reported in Section~\ref{sec:results} without re-simulating the working set separately per
capacity.

\section{Regularizing Distribution Shift Without a KL Penalty}
\label{app:sft}

A KL penalty against a frozen reference, or a distillation-style reward that scores a
rollout by the reference's own log-likelihood of it, both anchor the policy to the
reference's original behavior, but neither one ever increases the likelihood the policy
assigns to a completion known to be correct, since both are only ever satisfied by
staying close to what the reference already does. The supervised term in
Equation~\ref{eq:combinedloss} instead directly minimizes the negative log-likelihood of
the ground truth completion under the current policy, which pushes that likelihood up
regardless of what the reference would have assigned it, while the working-set reward
simultaneously pushes the policy toward whichever sampled rollout was most working-set
friendly. In our runs, we find the KL penalty to be redundant once the supervised term
is enabled, and disable it accordingly, relying on the supervised term alone to keep the
policy from drifting into working-set-friendly but degenerate completions.

\section{Hyperparameters}
\label{app:hparams}

\begin{table}[h]
\centering
\caption{Training hyperparameters shared across objectives, per backbone.}
\begin{tabular}{lcc}
\toprule
 & Phi-tiny-MoE-instruct & OLMoE-1B-7B \\
\midrule
Experts / top-$K$ & 16 / 2 & 64 / 8 \\
Layers & 32 & 16 \\
Working-set layer & 16 & 8 \\
LoRA rank / alpha & 16 / 32 & 16 / 32 \\
Group size $G$ & 8 & 8 \\
Learning rate & $1\mathrm{e}{-4}$ & $1\mathrm{e}{-4}$ \\
Prompt / completion length & 1024 / 1024 & 1024 / 1024 \\
\bottomrule
\end{tabular}
\end{table}

\end{document}
